Iloraz przepustowości rośnie z~liczbą klientów dla wszystkich par i~rozmiarów.
Współczynnik Spearmana wynosi wszędzie 1 (zależność ściśle monotoniczna). Dla
10 mln biletów przewaga Valkey nad PostgreSQL rośnie
z~\cratio{valkey}{postgres}{10000000}{lua}{cond}{1}$\times$ przy jednym kliencie
do \cratio{valkey}{postgres}{10000000}{lua}{cond}{1000}$\times$ przy tysiącu.
Nad MariaDB rośnie z~\cratio{valkey}{mariadb}{10000000}{lua}{cond}{1}$\times$
do \cratio{valkey}{mariadb}{10000000}{lua}{cond}{1000}$\times$.

Wynik dla MongoDB trzeba czytać ostrożniej. Iloraz MongoDB do PostgreSQL rośnie
z~\cratio{mongo}{postgres}{10000000}{atomic}{cond}{1} do
\cratio{mongo}{postgres}{10000000}{atomic}{cond}{1000}: przy małej
współbieżności MongoDB jest kilka razy wolniejszy, a~przy 1000 klientach
dorównuje PostgreSQL. Względna pozycja MongoDB rośnie, ale realną przewagę
ma on tylko nad MariaDB. Hipotezę H1 uznajemy za
\textbf{potwierdzoną dla Valkey}. Dla MongoDB potwierdza się teza
o~rosnącym stosunku wydajności, a~nie o~przewadze w~każdym punkcie.

Dwa czynniki wzmacniają wynik Valkey i~trzeba je brać pod uwagę: brak zapisu
na dysk oraz łączenie poleceń w~potoki przez sterownik. Strategia \texttt{watch}
pokazuje też, że ten sam silnik z~nieodpowiednią metodą synchronizacji
przegrywa ze wszystkimi bazami relacyjnymi. H1 jest więc prawdziwa dla
właściwie zaprojektowanej operacji atomowej, a~nie dla silnika
niezależnie od sposobu użycia.
